\section{Алгоритмы идентификации параметров}

\subsection{Vector Fitting}

Метод \textit{Vector Fitting} (VF) \cite{triverio2019vector} предназначен для построения редуцированных моделей линейных стационарных систем по дискретным выборкам их частотной характеристики. 
Пусть заданы $\bar{k}$ измерений передаточной функции:
\begin{equation}
    H_k = H(j\omega_k), \quad k=1,\dots,\bar{k}.
\end{equation}
Требуется найти рациональную аппроксимацию $\tilde{H}(s)$, минимизирующую среднеквадратичную ошибку:
\begin{equation}
    e^2 = \frac{1}{\bar{k}} \sum_{k=1}^{\bar{k}} \left| H_k - \tilde{H}(j\omega_k) \right|^2.
\end{equation}
Прямое решение этой нелинейной задачи наименьших квадратов (НКМ) затруднено из-за неизвестных коэффициентов в знаменателе, что приводит к плохой обусловленности (матрицы Вандермонда) и риску попадания в локальные минимумы.

VF заменяет полиномиальный базис $s^n$ на базис частичных дробей с фиксированными на текущей итерации полюсами $p_n^{(i-1)}$. Аппроксимант и весовая функция записываются в виде:
\begin{align}
    \tilde{H}^{(i)}(s) &= r_0^{(i)} + \sum_{n=1}^{\bar{n}} \frac{r_n^{(i)}}{s - p_n^{(i-1)}}, \\
    w^{(i)}(s) &= 1 + \sum_{n=1}^{\bar{n}} \frac{w_n^{(i)}}{s - p_n^{(i-1)}},
\end{align}
где $\bar{n}$ -- порядок модели. Вводится линеаризованный функционал ошибки:
\begin{equation}
    (e_{\text{SK}}^{(i)})^2 = \frac{1}{\bar{k}} \sum_{k=1}^{\bar{k}} \left| H_k w^{(i)}(j\omega_k) - \tilde{H}^{(i)}(j\omega_k) \right|^2.
\end{equation}
В отличие от классического метода Санатханана--Кёрнера, весовая функция $w^{(i)}(s)$ применяется \textit{неявно} через перенос полюсов на каждой итерации, что радикально улучшает обусловленность системы уравнений.

Алгоритм выполняется циклически ($i=1,2,\dots$):

\begin{enumerate}
    \item \textbf{Инициализация полюсов.} Выбираются начальные полюса $p_n^{(0)}$ с малой отрицательной вещественной частью, равномерно распределённые по диапазону частот $[\omega_1, \omega_{\bar{k}}]$:
    \begin{equation}
        p_n^{(0)} = \begin{cases}
            (-\alpha + j) \frac{\omega_{\bar{k}}}{\bar{n}/2} n, & n=1,\dots,\bar{n}/2,\\
            (p_{n-\bar{n}/2}^{(0)})^*, & n=\bar{n}/2+1,\dots,\bar{n},
        \end{cases}
    \end{equation}
    где обычно $\alpha = 0.01$.
    
    \item \textbf{Решение линейной системы НКМ.} Минимизация $(e_{\text{SK}}^{(i)})^2$ сводится к решению:
    \begin{equation}
        \begin{bmatrix}
            \bm{\Phi}_0^{(i)} & -\mathbf{D}_H \bm{\Phi}_1^{(i)}
        \end{bmatrix}
        \begin{bmatrix}
            \mathbf{c}_H^{(i)} \\ \mathbf{c}_w^{(i)}
        \end{bmatrix}
        = \mathbf{V}_H,
    \end{equation}
    где $\bm{\Phi}_0^{(i)}, \bm{\Phi}_1^{(i)}$ содержат значения базисных функций $\frac{1}{j\omega_k - p_n^{(i-1)}}$, $\mathbf{D}_H = \operatorname{diag}(H_1,\dots,H_{\bar{k}})$, $\mathbf{V}_H = [H_1,\dots,H_{\bar{k}}]^T$, а $\mathbf{c}_H^{(i)}, \mathbf{c}_w^{(i)}$ -- векторы неизвестных вычетов и весовых коэффициентов. Система решается через QR-разложение.
    
    \item \textbf{Обновление полюсов.} Новые полюса $p_n^{(i)}$ вычисляются как нули функции $w^{(i)}(s)$, что эквивалентно задаче на собственные значения:
    \begin{equation}
        \{p_n^{(i)}\} = \operatorname{eig}\left( \mathbf{A}^{(i-1)} - \mathbf{b}_w (\mathbf{c}_w^{(i)})^T \right),
    \end{equation}
    где $\mathbf{A}^{(i-1)} = \operatorname{diag}(p_1^{(i-1)},\dots,p_{\bar{n}}^{(i-1)})$, $\mathbf{b}_w = [1,\dots,1]^T$.
    
    \item \textbf{Проверка сходимости.} Итерации останавливаются, когда:
    \begin{equation}
        \max_k \left| w^{(i)}(j\omega_k) - 1 \right| \le \varepsilon_w.
    \end{equation}
    После выполнения условия фиксируются полюса $p_n^{(i)}$ и решается окончательная задача НКМ только для вычетов:
    \begin{equation}
        \bm{\Phi}_0^{(i+1)} \mathbf{c}_H^{(i+1)} = \mathbf{V}_H,
    \end{equation}
    что гарантирует минимизацию исходной ошибки $e^2$.
\end{enumerate}

Для систем с $\bar{m}$ входами и $\bar{q}$ выходами измеряемые данные представляют собой матрицы $\mathbf{H}_k \in \mathbb{C}^{\bar{q}\times\bar{m}}$. Модель строится с \textit{общими полюсами} для всех элементов передаточной матрицы:
\begin{equation}
    \tilde{\mathbf{H}}(s) = \mathbf{R}_0 + \sum_{n=1}^{\bar{n}} \frac{\mathbf{R}_n}{s - p_n},
\end{equation}
где $\mathbf{R}_n \in \mathbb{C}^{\bar{q}\times\bar{m}}$. Минимизируется норма Фробениуса ошибки. Для снижения вычислительной сложности применяется \textit{Fast VF}, использующий блочные QR-разложения и распараллеливание по каналам.

Для реальных систем полюса и вычеты должны быть вещественными или образовывать комплексно-сопряжённые пары. Это обеспечивается явным разделением базиса на вещественную и комплексную части. 

Чтобы гарантировать устойчивость и причинность модели, на шаге обновления полюсов применяется принудительная стабилизация:
\begin{equation}
    p_n^{(i)} \leftarrow \begin{cases}
        p_n^{(i)}, & \operatorname{Re}\{p_n^{(i)}\} < 0,\\
        -\operatorname{Re}\{p_n^{(i)}\} + j \operatorname{Im}\{p_n^{(i)}\}, & \operatorname{Re}\{p_n^{(i)}\} > 0.
    \end{cases}
\end{equation}
При отсутствии шума в данных это не снижает точность, но делает модель пригодной для временн\'{o}го моделирования.

Полученная модель в форме ``полюса--вычеты'' преобразуется к пространству состояний (реализация Гилберта):
\begin{equation}
    \begin{cases}
        \dot{\mathbf{x}}(t) = \mathbf{A}\mathbf{x}(t) + \mathbf{B}\mathbf{u}(t),\\
        \mathbf{y}(t) = \mathbf{C}\mathbf{x}(t) + \mathbf{D}\mathbf{u}(t),
    \end{cases}
\end{equation}
где $\mathbf{D}=\mathbf{R}_0$, а блочно-диагональная матрица $\mathbf{A}$ и блочные $\mathbf{B}, \mathbf{C}$ формируются из SVD-разложений матриц вычетов $\mathbf{R}_n = \mathbf{U}_n \bm{\Sigma}_n \mathbf{V}_n^H$. Такая реализация гарантирует минимальный порядок системы и исключает неуправляемые/ненаблюдаемые состояния.

Таким образом, метод Vector Fitting преобразует нелинейную задачу рациональной аппроксимации в последовательность линейных задач НКМ за счёт неявного весового перевзвешивания и переноса полюсов. Метод обладает высокой скоростью сходимости (обычно 4--6 итераций), численной устойчивостью и легко масштабируется на многоканальные системы. Существуют расширения алгоритма для работы во временной области ($z$-VF), учёта шумов (Relaxed VF), параметрического моделирования и распределённых систем с задержками.

\subsection{Алгоритмы на основе CNN-архитектур}

Электрохимическая импедансная спектроскопия (EIS) позволяет характеризовать электрохимические системы через измерение комплексного импеданса $Z$ в диапазоне частот $f = [f_1, f_2, \dots, f_M]$. Для интерпретации данных широко используется моделирование эквивалентными электрическими схемами (ECM), где импеданс описывается функцией от параметров компонентов:
\begin{equation}
    Z = F(f; p_1, p_2, \dots, p_N),
\end{equation}
где $\mathbf{p} = [p_1, \dots, p_N]$ --- вектор параметров (сопротивления, ёмкости, индуктивности, элементы постоянной фазы и т.д.).

Традиционный метод комплексных нелинейных наименьших квадратов (CNLS), реализуемый через алгоритм Левенберга--Марквардта, требует задания начального приближения $\mathbf{p}^{(0)}$. Сильная нелинейность модели и наличие шума часто приводят к сходимости в локальные минимумы или расходимости. Для решения проблемы авторы предлагают использовать сверточную нейронную сеть (CNN) для автоматической генерации качественного начального приближения \cite{hua2023deep}.

\begin{figure}[H]
	\centering
	\includegraphics [width=1.00\linewidth] {review/CNN_encoder.png}

	\caption{}
	\label{fig:CNN_encoder}
\end{figure}

Для обучения регрессионной модели генерируется синтетический набор данных. Параметры $\mathbf{p}$ выбираются случайно в физически обоснованном диапазоне, после чего с помощью симулятора рассчитываются соответствующие спектры импеданса. Входные данные нормализуются и представляются в виде 4 каналов: действительная часть ($Z'$), мнимая часть ($Z''$), модуль ($|Z|$) и фаза ($\angle Z$).

Архитектура сети включает:
\begin{itemize}[leftmargin=*]
    \item Три слоя \texttt{Conv1D} с количеством фильтров 128, 256, 128 и длинами ядер 4, 5, 4 соответственно.
    \item Слой \texttt{Flatten}.
    \item Слой \texttt{Gaussian Noise} со стандартным отклонением $\sigma = 0.05$ для повышения обобщающей способности.
    \item Три полносвязных (\texttt{Dense}) слоя размерности 256, 64 и $p$ (число параметров схемы).
\end{itemize}

Во всех слоях используется функция активации ReLU:
\begin{equation}
    \text{ReLU}(x) = \max(0, x).
\end{equation}

Обучение проводится с помощью оптимизатора Adam (learning rate $= 0.001$). Функция потерь определяется как среднеквадратичная ошибка (MSE) между предсказанными сетью параметрами $p'_n$ и истинными параметрами $p_n$:
\begin{equation}
    \mathcal{L}_{\text{CNN}} = \frac{1}{N} \sum_{n=1}^{N} (p'_n - p_n)^2.
\end{equation}

Обученная модель выступает в роли этапа предфиттинга. Процедура состоит из следующих шагов:
\begin{enumerate}[leftmargin=*]
    \item Измеренный спектр импеданса нормализуется и подаётся на вход CNN.
    \item Сеть генерирует начальное приближение параметров $\mathbf{p}_{\text{init}}$.
    \item $\mathbf{p}_{\text{init}}$ передаётся в алгоритм CNLS (Левенберг--Марквардт) для точной оптимизации.
\end{enumerate}

CNLS минимизирует ошибку между измеренными значениями $z_m$ и модельными $Z_m$ на каждой частоте:
\begin{equation}
    \text{MSE}_{\text{CNLS}} = \frac{1}{M} \sum_{m=1}^{M} (z_m - Z_m)^2.
\end{equation}
Благодаря качественному начальному приближению, алгоритм быстро сходится к глобальному минимуму без ручного вмешательства.

Для валидации метода используется средняя абсолютная процентная ошибка (MAPE), вычисляемая по всем частотным точкам:
\begin{equation}
    \text{MAPE} = \frac{1}{M} \sum_{m=1}^{M} \frac{|z_m - Z_m|}{|z_m|}.
\end{equation}
Фиттинг считается успешным, если $\text{MAPE} < 0.01$. Эксперименты на открытом датасете показали, что предложенный метод повышает успешность фиттинга до $87.5\%$ (на $25\%$ выше по сравнению с CNLS со случайными начальными значениями и стохастическими методами), при этом сокращая общее время вычислений и обеспечивая стабильную сходимость.

\subsection{Архитектуры для задач сегментации}

\begin{figure}[H]
	\centering
	\includegraphics [width=1.00\linewidth] {review/cracks_encoder_eecoder.png}

	\caption{Архитектура автоэнкодера для сегментации \enquote{узких} объектов}
	\label{fig:LPF_0}
\end{figure}

Метод \cite{kompanets2025deep} основан на архитектуре \textbf{кодировщик-декодировщик} (encoder-decoder), адаптированной для работы с высокоразрешающими изображениями (до $4608 \times 3456$ пикселей). Сеть выполняет бинарную посимвольную классификацию: каждому пикселю присваивается значение $0$ (фон) или $1$ (трещина).

В качестве экстрактора признаков используется архитектура \textbf{ConvNeXt} (версия \texttt{Tiny}), предобученная на наборе данных ImageNet. ConvNeXt состоит из 4 стадий, каждая из которых включает свёрточные блоки с увеличением глубины каналов и уменьшением пространственного разрешения. Использование предобученных весов реализует стратегию трансферного обучения, позволяя сети эффективно извлекать низкоуровневые и семантические признаки даже на ограниченном доменном наборе данных.

Декодировщик состоит из 6 уровней и модифицирует архитектуру \cite{konig2021optimized}. Основные компоненты:
\begin{itemize}
    \item \textbf{Модули внимания:} На каждом уровне применяются пространственные и канальные слои squeeze-and-excitation, которые адаптивно взвешивают карты признаков, усиливая релевантные области.
    \item \textbf{Skip-connections:} Карты признаков декодировщика конкатенируются с соответствующими картами кодировщика для восстановления пространственной информации, потерянной при даунсэмплинге.
    \item \textbf{Остаточные блоки:} После конкатенации и свёртки $3\times3$ применяются residual-блоки с batch normalization и функцией активации ReLU.
    \item \textbf{Финальный слой:} Поточечная свёртка с сигмоидальной активацией формирует итоговую карту сегментации со значениями в диапазоне $[0, 1]$.
\end{itemize}

Прямая подача полного изображения в сеть требует $\sim 189$ ГБ ОЗУ, что превышает возможности стандартных GPU (40 ГБ). Для решения проблемы применяется \textbf{патч-ориентированный подход}:
\begin{itemize}
    \item Изображение разбивается на непересекающиеся патчи размером $512 \times 512$ пикселей.
    \item Каждый патч обрабатывается сетью независимо.
    \item Полученные карты сегментации склеиваются (stitching) для формирования глобальной карты всего изображения.
\end{itemize}
Эксперименты показали, что уменьшение размера патча (например, до $128 \times 128$) снижает точность, так как сеть теряет глобальный контекст, необходимый для отличия трещин от \textit{crack-like features} (стыки, тени, коррозия).

Стандартной функцией потерь для сегментации трещин является сумма бинарной кросс-энтропии (BCE) и Dice-потери, которая компенсирует сильный класс-имбаланс (трещины занимают малую площадь):
\begin{align}
    L_{\text{BCE}}(y, \hat{y}) &= -\frac{1}{N} \sum_{i=1}^{N} \left[ y_i \log \hat{y}_i + (1 - y_i) \log(1 - \hat{y}_i) \right], \\
    L_{\text{Dice}}(y, \hat{y}) &= 1 - \frac{\sum_{i=1}^{N} y_i \hat{y}_i + \epsilon}{\sum_{i=1}^{N} y_i + \sum_{i=1}^{N} \hat{y}_i + \epsilon}, \\
    L(y, \hat{y}) &= L_{\text{BCE}}(y, \hat{y}) + L_{\text{Dice}}(y, \hat{y}),
\end{align}
где $N$ --- число пикселей, $y_i \in [0,1]$ и $\hat{y}_i \in \{0,1\}$ --- значения $i$-го пикселя в эталонной и предсказанной картах соответственно, $\epsilon=1$.

\textbf{Проблема фона:} Для патчей без трещин $\sum_{i=1}^{N} y_i = 0$. В этом случае $L_{\text{Dice}} \approx 1$ при любом $\sum \hat{y}_i \gg \epsilon$, что делает функцию нечувствительной к ложноположительным срабатываниям на фоне.

\textbf{Решение (Loss Inversion):} Для фоновых патчей вводится операция инверсии меток:
\begin{equation}
    L_{\text{inversion}}(y, \hat{y}) = 
    \begin{cases}
        L(\text{inv}(y), \text{inv}(\hat{y})), & \text{если } \sum_{i=1}^{N} y_i = 0, \\
        L(y, \hat{y}), & \text{иначе},
    \end{cases}
\end{equation}
где $\text{inv}(z) = -z + 1$. При обучении на фоне сеть оптимизируется как «сегментатор фона», что позволяет корректно штрафовать ошибки и использовать в обучающей выборке до 90\% фоновых патчей без попадания в локальный минимум «all-white segmentation maps».

Для стабилизации градиентов и ускорения сходимости на промежуточных уровнях декодировщика (уровни 2--5) формируются дополнительные карты сегментации. Итоговая функция потерь вычисляется как взвешенная сумма:
\begin{equation}
    L_{\text{total}}(y, \hat{y}) = L_{\text{target}}(y, \hat{y}) + \sum_{i=0}^{n} w_i L_{\text{companion}}^i(y, \hat{y}_i),
\end{equation}
где $n$ -- число вспомогательных карт, $w_i = 1$, а $\hat{y}_i$ масштабируются до размера входного изображения методом ближайших соседей и проходят через сигмоиду.

\begin{itemize}
    \item \textbf{Оптимизатор:} ADAMW с $\lambda_{\text{initial}} = 0.001$ и weight decay $1 \cdot 10^{-5}$.
    \item \textbf{Экспоненциальное затухание LR:} $\lambda = \lambda_{\text{initial}} \cdot \gamma^e$, где $\gamma = 0.99$, $e$ --- номер эпохи.
    \item \textbf{Поэтапное затухание LR (ConvNeXt):} $\lambda_{\text{stage}} = \lambda \cdot k^{N+1-n}$, где $k=0.7$, $N=4$ (число стадий), $n$ --- номер стадии.
    \item \textbf{Аугментация:} Флипы, повороты $\{0^\circ, 90^\circ, 180^\circ, 270^\circ\}$, изменение яркости/контраста/насыщения (75--125\%), сдвиг оттенка ($\pm 10\%$), аддитивный шум, случайное масштабирование и зум.
    \item \textbf{Баланс патчей:} Оптимальное соотношение трещина/фон для стандартной потери $\approx 70/30$. При использовании $L_{\text{inversion}}$ возможно обучение на полном наборе ($\sim 6.5/93.5$).
\end{itemize}